% =========================================================================
% capitulos/02_marco_teorico.tex
% Capítulo II: Marco Teórico y Tecnológico
% =========================================================================

\section{Revisión Teórica}

La revisión teórica tiene por finalidad fundamentar científicamente y técnicamente el desarrollo del proyecto de prácticas preprofesionales, garantizando que las decisiones de ingeniería adoptadas se sustenten en literatura académica y estándares consolidados.

\subsection{Antecedentes}

En este apartado se deben recopilar investigaciones previas, tesis de grado o proyectos similares vinculados a la solución planteada. Se recomienda clasificar los antecedentes en el ámbito internacional, nacional y local.

Por ejemplo, según \citet{pressman2014}, la adopción sistemática de modelos de ciclo de vida en la ingeniería de software permite mitigar riesgos técnicos y optimizar el aprovechamiento de los recursos durante cada fase de implementación.

\subsection{Bases Teóricas y Científicas}

\subsubsection{Arquitectura de Software y Principios de Diseño}
Las decisiones arquitectónicas determinan la estabilidad a largo plazo de cualquier solución tecnológica. Como destaca \citet{martin2017clean}, una arquitectura desacoplada permite aislar las reglas de negocio esenciales de las dependencias volátiles de infraestructura, frameworks de interfaz y sistemas gestores de bases de datos.

Entre los principios rectores que deben guiar el desarrollo de software de calidad destacan los principios SOLID:
\begin{itemize}
  \item \textbf{Single Responsibility Principle (SRP):} Cada módulo o clase debe tener una única razón para cambiar.
  \item \textbf{Open/Closed Principle (OCP):} Las entidades de software deben estar abiertas a la extensión, pero cerradas a la modificación.
  \item \textbf{Liskov Substitution Principle (LSP):} Los subtipos deben ser sustituibles por sus tipos base sin alterar el funcionamiento del programa.
  \item \textbf{Interface Segregation Principle (ISP):} Ningún cliente debe depender de métodos que no utiliza.
  \item \textbf{Dependency Inversion Principle (DIP):} Los módulos de alto nivel no deben depender de los de bajo nivel; ambos deben depender de abstracciones.
\end{itemize}

\subsubsection{Metodologías Ágiles de Desarrollo}
El marco de trabajo Scrum \citep{schwaber2020scrum} se fundamenta en el empirismo y el control de procesos mediante iteraciones cortas denominadas \textit{Sprints}. Este enfoque favorece la adaptación continua a requerimientos emergentes y la entrega progresiva de incrementos funcionales de valor.

\subsection{Definición de Términos Básicos}
Presente aquí una lista de términos técnicos o glosario que clarifique conceptos específicos empleados a lo largo del informe:
\begin{itemize}
  \item \textbf{API (Application Programming Interface):} Conjunto de definiciones y protocolos utilizados para diseñar e integrar el software de las aplicaciones.
  \item \textbf{Frontend:} Capa de presentación visual e interacción directa con el usuario final.
  \item \textbf{Backend:} Capa lógica de procesamiento, acceso a datos y comunicación con servicios externos.
  \item \textbf{Refactorización:} Proceso de reestructuración interna del código fuente sin alterar su comportamiento observable \citep{fowler2018refactoring}.
\end{itemize}
